\section{The GHZ state from an upper-triangular tensor}%
\label{sec:asymex_ghz}

For $N>0$, the unnormalized Greenberger--Horne--Zeilinger (GHZ) state on
$(\C^2)^{\otimes N}$ is
\begin{align}
    \ket{\mathrm{GHZ}_N}
     & =\ket{0}^{\otimes N}+\ket{1}^{\otimes N}.
    \notag
\end{align}
Here $\ket{0},\ket{1}$ form the orthonormal local basis. We use the
product-state and GHZ conventions of Ref.~\cite[Appendix~A]{Cirac2021Matrix}.
For a tensor $T=\{T^i\}_{i=0}^1$ of bond dimension $D$, the periodic trace
boundary defines
\begin{align}
    V^{(N)}(T)_\sigma
     & =\tr(T^{\sigma(0)}\cdots T^{\sigma(N-1)}),
    \label{eq:asymex_periodic_amplitude}          \\
    \ket{V^{(N)}(T)}
     & =\sum_{\sigma:\Fin{N}\to\{0,1\}}
    V^{(N)}(T)_\sigma\ket{\sigma},
    \notag
\end{align}
where $\Fin{N}=\{0,\ldots,N-1\}$ for $N>0$, $\Fin{0}=\varnothing$,
and $\ket{\sigma}=\bigotimes_{n=0}^{N-1}\ket{\sigma(n)}$.
At $N=0$, the matrix product is $\Id_D$, so the unique coefficient is
$D$ and the physical space is $\C$.
For a word $w=(i_1,\ldots,i_{|w|})$, write
$T^w=T^{i_1}\cdots T^{i_{|w|}}$, with $T^\varnothing=\Id_D$.

The standard GHZ tensor has bond dimension two. The construction below
instead uses a nonminimal bond-dimension-four upper-triangular tensor
with the same positive-length states. Its two scalar diagonal sectors
generate the constant product configurations. These finite-length
identities make no claim about thermodynamic symmetry breaking.

\begin{definition}[The two sectors and the four-dimensional source]
    \label{def:asymex_ghz}
    \lean{MPSTensor.ghzC, MPSTensor.ghzB}
    \leanok
    \uses{def:mps_tensor}
    On the alphabet $\{0,1\}$, define the scalar tensors $C_0,C_1$ and
    the standard GHZ tensor $G$ by
    \begin{align}
        C_0^0 & =(1),        & C_0^1 & =(0), \notag \\
        C_1^0 & =(0),        & C_1^1 & =(1), \notag \\
        G^0   & =\diag(1,0), & G^1   & =\diag(0,1).
        \notag
    \end{align}
    Thus $G^i=C_0^i\oplus C_1^i$, and both scalar tensors are normal.
    The constructed source is the bond-dimension-four tensor
    \begin{align}
        B^0
         & =\begin{pmatrix}
                0 & 1 & 2 & 0 \\
                0 & 0 & 1 & 1 \\
                0 & 0 & 1 & 3 \\
                0 & 0 & 0 & 0
            \end{pmatrix},
         &
        B^1
         & =\begin{pmatrix}
                1 & 1 & 0 & 2 \\
                0 & 0 & 2 & 0 \\
                0 & 0 & 0 & 1 \\
                0 & 0 & 0 & 0
            \end{pmatrix}.
        \notag
    \end{align}
\end{definition}

\begin{center}
    % Formula: the periodic trace tr(B^w) of the bond-dimension-four
    % source of this section.
    % Source: Definition~\ref{def:asymex_ghz}.
    % Ink-to-index map: each node is the source tensor $B$; the upper
    % legs carry the letters of the word.
    % Contraction set: the horizontal virtual legs, closed around the
    % ring by the trace; no leg is left open.
    % Boundary signature: (virtual-west=0 | virtual-east=0 |
    % physical-up=3 | physical-down=0) for the finite schematic shown.
    \tenkzkernel
    \begin{tenkz}[cols=4, boundary=periodic, physical=up]
        \tn[label pos=135, ports={90:physical:$i_1$}]{B} &
        \tn[label pos=135]{B} & \tn[skin=dots]{} &
        \tn[label pos=135, ports={90:physical:$i_N$}]{B}
    \end{tenkz}
\end{center}

\begin{theorem}[GHZ amplitudes and periodic vectors]
    \label{thm:asymex_ghz_state}
    \lean{MPSTensor.ghzC_eq_ghzSectorTensor, MPSTensor.ghzC_mpv,
        MPSTensor.ghzTensor_mpv, MPSTensor.ghzB_mpv,
        MPSTensor.ghzB_sameMPV₂Pos_ghzTensor, MPSTensor.ghzB_mpvState_eq}
    \leanok
    \uses{def:asymex_ghz}
    For $x\in\{0,1\}$, $C_x$ is the standard scalar GHZ sector,
    with $C_x^i=(\delta_{i,x})$. For every $N\geq0$ and configuration
    $\sigma:\Fin{N}\to\{0,1\}$,
    \begin{align}
        V^{(N)}(C_x)_\sigma
         & =\begin{cases}
                1, & \sigma(n)=x\text{ for every }n\in\Fin{N}, \\
                0, & \text{otherwise},
            \end{cases}
        \notag                                             \\
        V^{(N)}(G)_\sigma
         & =V^{(N)}(C_0)_\sigma+V^{(N)}(C_1)_\sigma.
        \notag
    \end{align}
    For $N>0$, the source has the same amplitudes and vector:
    \begin{align}
        V^{(N)}(B)_\sigma
         & =V^{(N)}(G)_\sigma
        =\delta_{\sigma,0^N}+\delta_{\sigma,1^N},
        \label{eq:asymex_ghz_amplitude}            \\
        \ket{V^{(N)}(B)}
         & =\ket{V^{(N)}(G)}=\ket{\mathrm{GHZ}_N}.
        \notag
    \end{align}
    Here $x^N$ denotes the constant configuration with value $x$, and
    $\delta$ is the Kronecker delta.
\end{theorem}

\begin{proof}
    \leanok
    \uses{thm:asymex_ghz_compression}
    For a scalar tensor $T$ with $T^i=(t_i)$, induction on $N$ gives
    $V^{(N)}(T)_\sigma=\prod_{n=0}^{N-1}t_{\sigma(n)}$.
    Applying this to $C_x$ gives $\prod_n\delta_{\sigma(n),x}$.
    For $G$, induction on $N$ gives
    \begin{align}
        G^{\sigma(0)}\cdots G^{\sigma(N-1)}
         & =\diag\left(\prod_n\delta_{\sigma(n),0},
        \prod_n\delta_{\sigma(n),1}\right).
        \notag
    \end{align}
    Taking the trace proves the formula for $G$. For $N>0$, the
    compression datum in Theorem~\ref{thm:asymex_ghz_compression} gives
    $V^{(N)}(B)_\sigma=\sum_{x=0}^1V^{(N)}(C_x)_\sigma$.
    Equality of every coefficient gives equality of the vectors.
\end{proof}

\begin{theorem}[GHZ norm and the empty-length exception]
    \label{thm:asymex_ghz_norm}
    \lean{MPSTensor.ghzB_mpvState_norm_sq,
        MPSTensor.ghzB_not_sameMPV₂_ghzTensor}
    \leanok
    \uses{def:asymex_ghz}
    For every $N>0$,
    \begin{align}
        \left\|\ket{V^{(N)}(B)}\right\|^2 & =2.
        \notag
    \end{align}
    Thus the normalized vector is
    $2^{-1/2}(\ket{0}^{\otimes N}+\ket{1}^{\otimes N})$.
    The source and standard GHZ tensor do not generate equal vectors at
    all lengths: their empty-length coefficients are
    $V^{(0)}(B)_\varnothing=4$ and $V^{(0)}(G)_\varnothing=2$.
\end{theorem}

\begin{proof}
    \leanok
    \uses{thm:asymex_ghz_state}
    For $N>0$, the configurations $0^N$ and $1^N$ are distinct.
    By~(\ref{eq:asymex_ghz_amplitude}),
    \begin{align}
        \left\|\ket{V^{(N)}(B)}\right\|^2
         & =\sum_\sigma\left|V^{(N)}(B)_\sigma\right|^2
        =\sum_\sigma(\delta_{\sigma,0^N}+\delta_{\sigma,1^N})
        =2.
        \notag
    \end{align}
    At length zero, the traces are $\tr(\Id_4)=4$ and
    $\tr(\Id_2)=2$, which rules out equality at every length.
\end{proof}

\begin{theorem}[Compression of the two sectors]
    \label{thm:asymex_ghz_compression}
    % Principal declaration: MPSTensor.ghzSectors_isReduction.
    \lean{MPSTensor.ghzSectors_compression, MPSTensor.ghzSectors_isReduction,
        MPSTensor.ghzSectors_dim_eq}
    \leanok
    \uses{def:asymex_ghz, thm:asym_multiblock}
    The source of Definition~\ref{def:asymex_ghz} carries the block
    triangular form of Theorem~\ref{thm:asym_multiblock} for the slot family
    $(C_0,C_1)$ with $z=2$ zero slots, the four diagonal positions carrying,
    in order, the sector $C_1$, a zero slot, the sector $C_0$ and a second
    zero slot. The compression pair of each sector satisfies $W_sV_s=\id$
    and $W_sB^wV_s=C_s^w$ for every word, and $4=1+1+2$.
    Here $V_s:\C\to\C^4$ and $W_s:\C^4\to\C$ are the coordinate
    inclusion and projection at position $2$ for $s=0$ and position $0$
    for $s=1$, with coordinates numbered from zero.
\end{theorem}

\begin{proof}
    \leanok
    \uses{thm:asymex_explicit_gauge}
    The gauge is the relabelling of the four bond coordinates by the four
    slots, with no change of basis. The three matrix conditions of the
    compression datum, block triangularity and the identification of the
    matched and unmatched diagonal blocks, follow by checking the entries
    of $B^0$ and $B^1$. The compression pairs and
    the dimension count are clauses (iv), (v) and (vii) of
    Theorem~\ref{thm:asym_multiblock}.
\end{proof}

\begin{theorem}[Periodic traces of the source]
    \label{thm:asymex_ghz_trace}
    \lean{MPSTensor.ghzB_trace_evalWord_eq_sum}
    \leanok
    \uses{def:asymex_ghz}
    For every nonempty word $w$, $\tr(B^w)=\tr(C_0^w)+\tr(C_1^w)$.
\end{theorem}

\begin{proof}
    \leanok
    \uses{thm:asymex_ghz_compression, prop:asym_compression_trace}
    This is Proposition~\ref{prop:asym_compression_trace} applied to the
    compression datum of Theorem~\ref{thm:asymex_ghz_compression}, whose
    slot set has two members.
\end{proof}

\begin{theorem}[The first sector has no sitewise intertwiner]
    \label{thm:asymex_ghz_no_intertwiner}
    \lean{MPSTensor.ghzC0_right_intertwiner_eq_zero,
        MPSTensor.ghzC0_left_intertwiner_eq_zero}
    \leanok
    \uses{def:asymex_ghz}
    If a column vector $v\in\C^4$ satisfies $B^0v=v$ and $B^1v=0$,
    then $v=0$; if a row vector $u\in\C^4$ satisfies $uB^0=u$ and
    $uB^1=0$, then $u=0$. Both sitewise intertwiner spaces for the sector
    $C_0$ therefore vanish. In particular, the word-level compression of
    Theorem~\ref{thm:asymex_ghz_compression} does not give a nonzero
    sitewise intertwiner in either direction.
\end{theorem}

\begin{proof}
    \leanok
    The last row of $B^0$ vanishes, so the last coordinate of $B^0v=v$
    gives $v_3=0$; the second coordinate of $B^1v=0$ gives $v_2=0$; and the
    second and first coordinates of $B^0v=v$ then force $v_1=0$ and
    $v_0=0$. For the covector, the first and third coordinates of $uB^1=0$
    give $u_0=0$ and $u_1=0$, its last coordinate gives $u_2=0$, and the
    last coordinate of $uB^0=u$ gives $u_3=0$.
\end{proof}

\section{A product state from a repeated block}%
\label{sec:asymex_repeated}

Let $\ket{\phi}=\ket{0}+2\ket{1}$. The unnormalized product vector
$\ket{\phi}^{\otimes N}$ has a bond-one presentation with local
amplitudes $(1,2)$, as in the product-state construction of
Ref.~\cite[Appendix~A]{Cirac2021Matrix}. Repeating that scalar block twice
multiplies every periodic coefficient by two. The upper-triangular
source below generates the same vector as this repeated target:
\begin{align}
    \ket{V^{(N)}(B)}
     & =2\ket{\phi}^{\otimes N}
    =2(\ket{0}+2\ket{1})^{\otimes N}.
    \notag
\end{align}
In this section, $A$ and $B$ denote the tensors defined below, and
the trace boundary is again~(\ref{eq:asymex_periodic_amplitude}).
The tensor power at $N=0$ is the scalar $1$.

\begin{definition}[The scalar block and its Jordan extension]
    \label{def:asymex_repeated}
    \lean{MPSTensor.repA, MPSTensor.repC, MPSTensor.repB}
    \leanok
    \uses{def:mps_tensor}
    The target is the scalar normal tensor $A^0=(1)$, $A^1=(2)$, taken
    twice as $C_s=A$ for $s\in\{0,1\}$, and the source is
    \begin{align}
        B^0
         & =\begin{pmatrix}1&1\\0&1\end{pmatrix},
         &
        B^1
         & =\begin{pmatrix}2&0\\0&2\end{pmatrix}.
        \notag
    \end{align}
\end{definition}

\begin{center}
    % Formula: the periodic trace tr(B^w) of the bond-dimension-two
    % Jordan source of this section.
    % Source: Definition~\ref{def:asymex_repeated}.
    % Ink-to-index map: each node is the source tensor $B$; the upper
    % legs carry the letters of the word.
    % Contraction set: the horizontal virtual legs, closed around the
    % ring by the trace; no leg is left open.
    % Boundary signature: (virtual-west=0 | virtual-east=0 |
    % physical-up=3 | physical-down=0) for the finite schematic shown.
    \tenkzkernel
    \begin{tenkz}[cols=4, boundary=periodic, physical=up]
        \tn[label pos=135, ports={90:physical:$i_1$}]{B} &
        \tn[label pos=135]{B} & \tn[skin=dots]{} &
        \tn[label pos=135, ports={90:physical:$i_N$}]{B}
    \end{tenkz}
\end{center}

\begin{theorem}[Product-state amplitudes and multiplicity]
    \label{thm:asymex_repeated_state}
    \lean{MPSTensor.repA_mpv, MPSTensor.repB_mpv_eq_two_mul,
        MPSTensor.repB_mpv, MPSTensor.repB_mpvState_eq}
    \leanok
    \uses{def:asymex_repeated}
    Set $a_0=1$ and $a_1=2$. For every $N\geq0$ and
    $\sigma:\Fin{N}\to\{0,1\}$,
    \begin{align}
        V^{(N)}(A)_\sigma
         & =\prod_{n=0}^{N-1}a_{\sigma(n)}, \notag \\
        V^{(N)}(B)_\sigma
         & =2V^{(N)}(A)_\sigma
        =2\prod_{n=0}^{N-1}a_{\sigma(n)}.
        \label{eq:asymex_repeated_amplitude}
    \end{align}
    Consequently,
    \begin{align}
        \ket{V^{(N)}(A)}
         & =(\ket{0}+2\ket{1})^{\otimes N}, \notag \\
        \ket{V^{(N)}(B)}
         & =2\ket{V^{(N)}(A)}.
        \notag
    \end{align}
    The factor two also holds at $N=0$, where the coefficients of $A$
    and $B$ are $1$ and $2$, respectively.
\end{theorem}

\begin{proof}
    \leanok
    \uses{thm:asymex_repeated_compression}
    Induction on $N$ expresses the coefficient of the scalar tensor $A$
    as the product of its scalar entries. Substituting $A^0=(1)$ and
    $A^1=(2)$ gives the local amplitudes $a_0,a_1$. If $N>0$, the trace
    identity in Theorem~\ref{thm:asymex_repeated_compression} gives
    $V^{(N)}(B)_\sigma=2V^{(N)}(A)_\sigma$.
    For $N=0$, this follows from $\tr(\Id_2)=2\tr(\Id_1)$.
    Equality of the coefficients gives the vector identity.
\end{proof}

\begin{theorem}[Product-state norms]
    \label{thm:asymex_repeated_norm}
    \lean{MPSTensor.repA_mpvState_norm_sq, MPSTensor.repB_mpvState_norm_sq}
    \leanok
    \uses{def:asymex_repeated}
    For every $N\geq0$,
    \begin{align}
        \left\|\ket{V^{(N)}(A)}\right\|^2 & =5^N, \notag \\
        \left\|\ket{V^{(N)}(B)}\right\|^2 & =4\,5^N.
        \notag
    \end{align}
    Both tensors therefore give the same normalized product vector,
    $((\ket{0}+2\ket{1})/\sqrt{5})^{\otimes N}$.
\end{theorem}

\begin{proof}
    \leanok
    \uses{thm:asymex_repeated_state}
    Using the amplitudes $a_0=1$ and $a_1=2$
    in~(\ref{eq:asymex_repeated_amplitude}), expand the squared norm
    in the configuration basis and factor the finite sum:
    \begin{align}
        \left\|\ket{V^{(N)}(A)}\right\|^2
         & =\sum_\sigma\prod_{n=0}^{N-1}|a_{\sigma(n)}|^2
        =\prod_{n=0}^{N-1}(|a_0|^2+|a_1|^2)
        =5^N.
        \notag
    \end{align}
    The vector identity in Theorem~\ref{thm:asymex_repeated_state}
    multiplies the squared norm by $|2|^2=4$.
\end{proof}

The source and the direct sum $A\oplus A$ have the same periodic
vectors, including at length zero. Their equality does not imply an
invertible gauge between the tensors: the off-diagonal entry of $B^0$
does not contribute to any periodic word trace, but prevents conjugation
to the repeated target.

\begin{theorem}[Compression onto the two copies]
    \label{thm:asymex_repeated_compression}
    % Principal declaration: MPSTensor.repeatedBlock_isReduction.
    \lean{MPSTensor.repeatedBlock_compression, MPSTensor.repeatedBlock_isReduction,
        MPSTensor.repB_trace_evalWord_eq_two_mul}
    \leanok
    \uses{def:asymex_repeated, thm:asym_multiblock}
    The source of Definition~\ref{def:asymex_repeated} carries the block
    triangular form of Theorem~\ref{thm:asym_multiblock} for the two copies
    of $A$ with no zero slots. For $s\in\{0,1\}$, let $V_s:\C\to\C^2$
    and $W_s:\C^2\to\C$ be the inclusion and projection of coordinate $s$.
    Each coordinate pair compresses its copy:
    $W_sV_s=\id$ and $W_sB^wV_s=A^w$ for every word.
    Also, $\tr(B^w)=2\tr(A^w)$ for every nonempty word.
\end{theorem}

\begin{proof}
    \leanok
    \uses{prop:asym_compression_trace}
    Both matrices are upper triangular with both diagonal entries equal to
    the corresponding scalar of $A$, which verifies the three
    matrix conditions. The compression pairs are clauses (iv) and (v) of
    Theorem~\ref{thm:asym_multiblock}, and the trace identity is
    Proposition~\ref{prop:asym_compression_trace} with two equal slots.
\end{proof}

\begin{theorem}[No invertible gauge onto the repeated target]
    \label{thm:asymex_repeated_not_gauge}
    \lean{MPSTensor.not_gaugeEquiv_directSum_repA}
    \leanok
    \uses{def:asymex_repeated, def:gauge_equiv}
    There is no invertible $X\in\GL_2(\C)$ with
    $B^i=X\diag(A^i,A^i)X^{-1}$ for both letters $i$.
    Here $\diag(A^i,A^i)=A^i\oplus A^i$ is the repeated target matrix.
\end{theorem}

\begin{proof}
    \leanok
    For $i=0$ the matrix $\diag(A^0,A^0)$ is the identity, so any such
    conjugation would give $B^0=\Id_2$, contradicting the nonzero
    off-diagonal entry of $B^0$.
\end{proof}
